\section*{Capítulo \ref{ch:b2:aromatic-substitution} --- Aromaticidad y sustitución electrófila aromática}

\begin{solution}{exo:b2:aromatic-substitution:1}
\ce{HNO3 + 2H2SO4 -> NO2+ + H3O+ + 2HSO4-}: el ácido nítrico se protona y luego pierde agua.
\end{solution}

\begin{solution}{exo:b2:aromatic-substitution:2}
El tolueno da 2-bromotolueno y 4-bromotolueno: el metilo orienta a las posiciones orto y
para. El nitrobenceno da 3-bromonitrobenceno, más despacio: el grupo nitro
orienta a meta.
\end{solution}

\begin{solution}{exo:b2:aromatic-substitution:3}
\ce{-CH3}: activante, o/p. \ce{-OH}: fuertemente activante, o/p. \ce{-Cl}: desactivante, o/p.
\ce{-NO2}: fuertemente desactivante, meta. \ce{-COCH3}: desactivante, meta. \ce{-NHCOCH3}:
activante (moderadamente), o/p.
\end{solution}

\begin{solution}{exo:b2:aromatic-substitution:4}
Acetofenona (1-feniletan-1-ona), \ce{C6H5COCH3}, a través del \omterm{def:b2:aromatic-substitution:friedel-crafts}{ion acilio} \ce{CH3CO+}.
\end{solution}

\begin{solution}{exo:b2:aromatic-substitution:5}
El carbocatión primario (o su complejo) se transpone por un desplazamiento de hidruro al catión
isopropilo secundario, que alquila el anillo. Propilbenceno: \omterm{def:b2:aromatic-substitution:friedel-crafts}{acilación de Friedel--Crafts} con
cloruro de propanoílo (sin transposición) y luego reducción del \ce{C=O} de la cetona a \ce{CH2}.
\end{solution}

\begin{solution}{exo:b2:aromatic-substitution:6}
El grupo \ce{-OH} es un fuerte dador mesómero: el anillo está tan activado que el bromo molecular
basta como electrófilo, y se sustituyen las tres posiciones orto y para.
\end{solution}

\begin{solution}{exo:b2:aromatic-substitution:7}
3-Bromonitrobenceno: nítrese y luego brómese (el nitro orienta a meta). 4-Bromonitrobenceno:
brómese y luego nítrese (el bromo orienta a orto/para), y sepárese el isómero para del
orto.
\end{solution}

\begin{solution}{exo:b2:aromatic-substitution:8}
La mezcla nitrante oxida la anilina, muy rica en electrones, y protona su nitrógeno:
\ce{-NH3+} es un \omterm{def:b2:aromatic-substitution:directing}{grupo desactivante} y \omterm{def:b2:aromatic-substitution:directing}{orientador meta}. En la acetanilida el par no enlazante del nitrógeno está
compartido con el carbonilo: menos básico y menos activante, no se protona ni se oxida,
y sigue orientando a orto/para, principalmente a para por razones estéricas.
\end{solution}

\begin{solution}{exo:b2:aromatic-substitution:9}
Sulfónese el tolueno (sobre todo en para: el grupo voluminoso evita la posición orto); brómese: el bromo entra
en orto respecto al metilo (la posición para está bloqueada, y ambos grupos favorecen esa posición);
elimínese el grupo \ce{SO3H} calentando en ácido diluido.
\end{solution}

\begin{solution}{exo:b2:aromatic-substitution:10}
Ambos grupos orientan a orto/para; sus posiciones para están ocupadas por el otro grupo. El metoxi, el
activante más fuerte, gana: nitración en orto respecto al \ce{-OCH3} (posición 2), que da
4-metil-2-nitroanisol.
\end{solution}

\begin{solution}{exo:b2:aromatic-substitution:11}
Cada grupo nitro desactiva fuertemente el anillo: la segunda nitración requiere condiciones más fuertes
que la primera, y la tercera (sobre un anillo con dos grupos nitro y solo el metilo, débilmente
activante) más fuertes todavía.
\end{solution}

\begin{solution}{exo:b2:aromatic-substitution:12}
Ácido 4-nitrobenzoico: nítrese el tolueno, sepárese el isómero para y oxídese el metilo a
\ce{-COOH} (permanganato en caliente). Ácido 3-nitrobenzoico: oxídese primero el tolueno a ácido benzoico y luego
nítrese (el \ce{-COOH} orienta a meta).
\end{solution}

\begin{solution}{pb:b2:aromatic-substitution:1}
\textbf{1.} \ce{HNO3 + 2H2SO4 -> NO2+ + H3O+ + 2HSO4-}.
\textbf{2.} \ce{O=N+=O}: lineal, isoelectrónico con el \ce{CO2}.
\textbf{3.} El nitrógeno se une al carbono del anillo en para respecto al \ce{CH3}, que pasa a ser tetraédrico
(con H y \ce{NO2}).
\textbf{4.} La carga positiva sobre los dos carbonos en orto respecto al carbono atacado y sobre el
carbono en para respecto a él, que lleva el metilo.
\textbf{5.} La primera, la formación del \omterm{def:b2:aromatic-substitution:seAr}{intermedio de Wheland}.
\textbf{6.} Una base del medio: \ce{HSO4-} o el agua.
\textbf{7.} Dos orto, dos meta, una para.
\textbf{8.} 2 : 2 : 1, es decir, 40~\% orto, 40~\% meta, 20~\% para.
\textbf{9.} En el ataque en meta, la carga positiva nunca está sobre el carbono que lleva el metilo; ninguna
estructura recibe ayuda del metilo.
\textbf{10.} Una estructura de resonancia tiene la carga sobre el carbono que lleva el \ce{CH3}: un
carbocatión terciario, estabilizado por el efecto inductivo y la \omterm{def:b2:fragment-orbitals:hyperconjugation}{hiperconjugación} del metilo.
\textbf{11.} Orto y para están favorecidas (96~\% frente al 60~\% al azar); meta está casi ausente.
\textbf{12.} Por posición: orto $58/2 = 29~\%$, para 38~\%: la para es la posición más reactiva,
pues las posiciones orto están algo impedidas por el metilo.
\textbf{13.} Más deprisa: el metilo activa el anillo.
\textbf{14.} El 4-nitrotolueno (funde a \qty{52}{\degreeCelsius}).
\textbf{15.} Enfríese la mezcla bruta: el isómero para cristaliza y se separa por filtración; el
líquido conserva los isómeros orto y meta.
\textbf{16.} Una molécula simétrica se empaqueta mejor en un cristal: más interacciones por unidad de volumen,
un punto de fusión más alto.
\textbf{17.} Por \omterm{def:b2:liquid-vapour-diagrams:fractional-distillation}{destilación fraccionada} a presión reducida, o por cromatografía; o por un
enfriamiento adicional (el isómero meta funde a \qty{16}{\degreeCelsius}).
\textbf{18.} Los isómeros tienen puntos de ebullición próximos (misma fórmula, polaridad parecida).
\textbf{19.} \ce{C7H8}: \qty{92.14}{g/mol}; \ce{C7H7NO2}: \qty{137.14}{g/mol}.
\textbf{20.} $100/92.14 = \qty{1.085}{mol}$ de tolueno; $0.90 \times 1.085 = \qty{0.977}{mol}$ de
nitrotoluenos.
\textbf{21.} Total $0.977 \times 137.14 = \qty{134.0}{g}$: orto \qty{77.7}{g}, meta
\qty{5.4}{g}, para \qty{50.9}{g}.
\textbf{22.} Una segunda nitración (dinitrotoluenos), y la oxidación del grupo metilo.
\textbf{23.} Oxídese el metilo a \ce{-COOH} con permanganato alcalino en caliente y luego acidúlese.
\textbf{24.} La operación da $\boldsymbol{\approx \qty{51}{g}}$ de 4-nitrotolueno.
\end{solution}
